% Part: sets-functions-relations
% Chapter: size-of-sets
% Section: enumerations
% Hindi translation (hi-Deva-IN) of Open Logic commit 9620cc73f9c8e0ad003c514a5d3748f29611c4c0.

\documentclass[../../../include/open-logic-section]{subfiles}

\begin{document}

\olfileid[hi]{sfr}{siz}{enm}

\olsection{गणन और \usetoken{S}{enumerable} समुच्चय}

\begin{editorial}
  यह खंड समुच्चयों के गणन पर चर्चा करता है और उन्हें $\PosInt$ से आने वाले
  आच्छादक फलनों के रूप में परिभाषित करता है। कम गणितीय पृष्ठभूमि वाले पाठकों
  के लिए इसमें विषय को धीरे-धीरे समझाया गया है। एक वैकल्पिक, अधिक संक्षिप्त
  रूप \olref[enm-alt]{sec} में दिया गया है, जो गणन को अलग ढंग से परिभाषित
  करता है: $\Nat$ (या उसके किसी आरंभिक खंड) के साथ एकैकी आच्छादक प्रतिचित्रण
  के रूप में।
\end{editorial}

\begin{explain}
हम समुच्चयों के !!{element}s सूचीबद्ध करके उनके उदाहरण पहले ही दे चुके हैं।
अब अधिक सामान्य रूप में चर्चा करें कि किसी समुच्चय का प्रत्येक !!{element}s
कैसे और कब सूचीबद्ध किया जा सकता है, भले ही वह समुच्चय अनंत हो।
\end{explain}

\begin{defn}[गणन, अनौपचारिक रूप में]
अनौपचारिक रूप से, समुच्चय~$A$ का \emph{गणन} ऐसी (संभवतः अनंत) सूची है
जिसमें~$A$ के !!{element}s होते हैं और $A$ का प्रत्येक !!{element} किसी परिमित स्थान
पर आता है। यदि $A$ का कोई गणन है, तो $A$ को \emph{!!{enumerable}} कहते हैं।
\end{defn}

\begin{explain}
गणन के बारे में कुछ बातें ध्यान देने योग्य हैं:
\begin{enumerate}
\item हम केवल उन्हीं सूचियों को गणन मानते हैं जिनका आरंभ हो और जिनमें प्रथम
  को छोड़कर प्रत्येक !!{element} से ठीक पहले एक ही !!{element} हो। दूसरे
  शब्दों में, सूची के प्रथम !!{element} और किसी भी अन्य !!{element}
  के बीच केवल परिमित संख्या में अवयव होते हैं। विशेषतः इसका अर्थ है कि गणन के
  प्रत्येक !!{element} का स्थान परिमित होता है: प्रथम !!{element} का
  स्थान~$1$, दूसरे का स्थान~$2$, इत्यादि।
\item एक ही समुच्चय~$A$ के अलग-अलग गणन हो सकते हैं, जिनमें !!{element}s के आने
  का क्रम भिन्न हो: $4$, $1$, $25$, $16$,~$9$ प्रथम पाँच वर्ग संख्याओं (के
  समुच्चय) का उतना ही अच्छा गणन है जितना $1$, $4$, $9$, $16$,~$25$।
\item पुनरुक्तिपूर्ण गणन भी गणन ही होते हैं: $1$, $1$, $2$, $2$, $3$,
  $3$,~\dots{} उसी समुच्चय का गणन करता है जिसका $1$, $2$, $3$,~\dots{} करता
  है।
\item किसी गणन को निर्दिष्ट करते समय क्रम और पुनरुक्ति का \emph{वास्तव में}
  महत्त्व होता है: हम धन पूर्णांकों का गणन $1$, $2$, $3$, $1$, \dots{} से
  आरंभ कर सकते हैं, किंतु मानक ढंग से $1$, $2$, $3$, $4$,~\dots{} लिखने पर
  प्रतिरूप अधिक आसानी से दिखता है।
\item गणन का आरंभ होना आवश्यक है: \dots, $3$, $2$, $1$ धन पूर्णांकों का गणन
  नहीं है, क्योंकि उसका कोई प्रथम !!{element} नहीं है। यह अनौपचारिक परिभाषा
  से कैसे निकलता है, इसे देखने के लिए स्वयं से पूछें, ``सूची में संख्या 76 किस
  स्थान पर आती है?''
\item निम्नलिखित धन पूर्णांकों का गणन नहीं है: $1$, $3$, $5$, \dots, $2$,
  $4$, $6$, \dots\@ समस्या यह है कि सम संख्याएँ परिमित स्थानों के बजाय
  $\infty + 1$, $\infty + 2$, $\infty +
  3$ स्थानों पर आती हैं।
\item रिक्त समुच्चय गणनीय है: उसका गणन रिक्त सूची करती है!{}
\end{enumerate}
\end{explain}

\begin{prop}
  यदि $A$ का कोई गणन है, तो उसका पुनरावृत्तियों-रहित गणन भी है।
\end{prop}

\begin{proof}
  मान लें कि $A$ का एक गणन $x_1$, $x_2$, \dots{} है, जिसमें प्रत्येक $x_i$,
  ~$A$ का !!{element} है। दोहराए गए !!{element}s हटाकर हम गणन से
  पुनरावृत्तियाँ हटा सकते हैं। उदाहरणार्थ, हम गणन को ऐसी नई सूची में बदल सकते
  हैं जिसमें $x_i$ को तब लिखें जब वह~$A$ का ऐसा !!a{element} हो जो $x_1$,
  \dots, $x_{i-1}$ में न हो, और यदि $x_i$ पहले ही $x_1$,
  \dots,~$x_{i-1}$ में आ चुका हो तो उसे सूची से हटा दें।
\end{proof}

पिछली युक्ति दिखाती है कि गणन और !!{enumerable} समुच्चयों को अच्छी तरह समझने
और उनके विषय में कथन सिद्ध करने के लिए हमें अधिक सटीक परिभाषा चाहिए। वह
परिभाषा निम्नलिखित है।

\begin{defn}[गणन, औपचारिक रूप में] 
समुच्चय $A \neq \emptyset$ का \emph{गणन}, कोई भी !!{surjective} फलन
$f \colon \PosInt \to A$ है।
\end{defn}

\begin{explain}
आइए स्वयं को विश्वास दिलाएँ कि औपचारिक परिभाषा और संभवतः अनंत सूची वाली
अनौपचारिक परिभाषा समतुल्य हैं। पहले, $\PosInt$ से समुच्चय~$A$ में जाने वाला
कोई भी !!{surjective} फलन~$A$ का गणन करता है। ऐसा फलन ऊपर दी गई अनौपचारिक
परिभाषा के अनुसार एक गणन निर्धारित करता है: सूची $f(1)$, $f(2)$, $f(3)$,
\dots। क्योंकि $f$ !!{surjective} है,~$A$ का प्रत्येक !!{element},~$f(n)$
का मान अवश्य है, जहाँ~$n \in \PosInt$। अतः $A$ का प्रत्येक
!!{element} सूची में किसी परिमित स्थान पर आता है। चूँकि फलन का
!!{injective} होना आवश्यक नहीं, सूची पुनरुक्तिपूर्ण हो सकती है; किंतु यह
स्वीकार्य है (जैसा ऊपर बताया गया है)।

दूसरी ओर, ऐसी सूची दी हो जिसमें~$A$ के सभी !!{element}s आते हों, तो हम
!!a{surjective} फलन $f\colon \PosInt \to A$ इस प्रकार परिभाषित कर सकते हैं:
$f(n)$ को सूची का $n$वाँ !!{element} मानें, अथवा यदि $n$वाँ !!{element} न
हो तो सूची के अंतिम !!{element} को उसका मान लें। इससे !!a{surjective} फलन
केवल तब नहीं बनता जब $A$ रिक्त हो और इसलिए सूची भी रिक्त हो। अतः प्रत्येक
अरिक्त सूची कोई !!a{surjective} फलन $f\colon \PosInt \to A$ निर्धारित
करती है।
\end{explain}

\begin{defn}
  \ollabel{defn:enumerable}
  समुच्चय~$A$ !!{enumerable} है यदि और केवल यदि वह रिक्त है अथवा उसका कोई
  गणन है।
\end{defn}

\begin{ex}
धन पूर्णांकों ($\PosInt$) का गणन करने वाला फलन केवल $f(n) = n$ से दिया गया
तत्समक फलन है। प्राकृत संख्याओं $\Nat$ का गणन करने वाला फलन
$g(n) = n - 1$ है।
\end{ex}

\begin{ex}
$f\colon \PosInt \to \PosInt$ और $g \colon \PosInt \to
\PosInt$ फलन, जो
\begin{align*}
f(n) & = 2n \text{ तथा}\\
g(n) & = 2n - 1
\end{align*}
से दिए गए हैं, क्रमशः सम धन पूर्णांकों और विषम धन पूर्णांकों का गणन करते
हैं। किंतु इनमें से कोई भी फलन $\PosInt$ का गणन नहीं है, क्योंकि दोनों में
से कोई !!{surjective} नहीं है।
\end{ex}

\begin{prob}
धन वर्गों $1$, $4$, $9$, $16$, \dots{} का एक गणन परिभाषित कीजिए।
\end{prob}

\begin{ex}
फलन $f(n) = (-1)^{n} \lceil \frac{(n-1)}{2}\rceil$ (जहाँ
$\lceil x \rceil$, \emph{लघुतम पूर्णांक फलन} (छत फलन) को दर्शाता है, जो
$x$ को उससे छोटी न होने वाली निकटतम पूर्णांक संख्या तक पूर्णांकित करता है)
पूर्णांकों के समुच्चय~$\Int$ का गणन करता है। ध्यान दें कि $f$, धन और ऋण
पूर्णांकों के बीच आगे-पीछे ``उछलते'' हुए $\Int$ के मान कैसे उत्पन्न करता है:
\[
\begin{array}{c c c c c c c c}
f(1) & f(2) & f(3) & f(4) & f(5) & f(6) & f(7) & \dots \\ \\
- \lceil \tfrac{0}{2} \rceil & \lceil \tfrac{1}{2}\rceil & - \lceil \tfrac{2}{2} \rceil & \lceil \tfrac{3}{2} \rceil & - \lceil \tfrac{4}{2} \rceil  & \lceil \tfrac{5}{2}
\rceil & - \lceil \tfrac{6}{2} \rceil & \dots \\ \\
0 & 1 & -1 & 2 & -2 & 3 & \dots
\end{array}
\]
आप $f$ को निम्नलिखित प्रकार से अलग-अलग स्थितियों में परिभाषित भी मान सकते हैं:
\[
f(n) = \begin{cases}
  0 & \text{यदि $n = 1$}\\
  n/2 & \text{यदि $n$ सम है}\\
  -(n-1)/2 & \text{यदि $n$ विषम है और $>1$ है}
  \end{cases}
\]
\end{ex}

\begin{prob}
  दिखाइए कि यदि $A$ और $B$ !!{enumerable} हैं, तो $A \cup B$ भी गणनीय है।
  इसके लिए मान लीजिए कि !!{surjective} फलन $f\colon \PosInt \to
  A$ और $g\colon \PosInt \to B$ हैं; फिर !!a{surjective}
  फलन~$h\colon \PosInt \to A \cup B$ परिभाषित करके सिद्ध कीजिए कि वह
  !!{surjective} है। उन स्थितियों पर भी विचार कीजिए जहाँ $A$ या~$B = \emptyset$।
\end{prob}
  
\begin{prob}
  दिखाइए कि यदि $B \subseteq A$ और $A$ !!{enumerable} है, तो~$B$ भी गणनीय
  है। इसके लिए मान लीजिए कि कोई !!a{surjective} फलन $f\colon \PosInt \to
  A$ है। कोई !!a{surjective} फलन~$g\colon \PosInt \to B$ परिभाषित कीजिए
  और सिद्ध कीजिए कि वह !!{surjective} है। यदि $B = \emptyset$ हो तो क्या
  होता है?
\end{prob}
    
\begin{prob}
  $n$ पर आगमन द्वारा दिखाइए कि यदि $A_1$, $A_2$, \dots, $A_n$ सभी
  !!{enumerable} हैं, तो $A_1 \cup \dots \cup A_n$ भी गणनीय है। आप इस
  तथ्य को मान सकते हैं कि यदि दो समुच्चय $A$ और~$B$ !!{enumerable} हैं, तो
  $A \cup
  B$ भी गणनीय है। 
\end{prob}

यद्यपि किसी समुच्चय का प्रत्येक !!{element}s सूचीबद्ध करते समय $1$वें !!{element}
से गिनना आरंभ करना अधिक स्वाभाविक लग सकता है, गणितज्ञ वस्तुओं को गिनने के लिए
प्राकृत संख्याओं~$\Nat$ का उपयोग करना पसंद करते हैं। वे सूची के $0$वें,
$1$वें, $2$वें इत्यादि स्थान पर आने वाले प्रत्येक !!{element}s की बात करते हैं। तदनुसार, हम गणन को
$\Nat$ से~$A$ में जाने वाले !!a{surjective} फलन के रूप में परिभाषित कर
सकते हैं। निस्संदेह दोनों परिभाषाएँ समतुल्य हैं।

\begin{prop}\ollabel{prop:enum-shift}
  कोई !!a{surjection} $f\colon \PosInt \to A$ है यदि और केवल यदि कोई
  !!a{surjection} $g\colon \Nat \to A$ है।
\end{prop}

\begin{proof}
  कोई !!a{surjection} $f\colon \PosInt \to A$ दिया हो, तो हम
  $g(n) =
  f(n+1)$ को सभी $n \in \Nat$ के लिए परिभाषित कर सकते हैं। यह देखना सरल
  है कि $g\colon \Nat
  \to A$ !!{surjective} है। इसके विपरीत, कोई
  !!a{surjection} $g\colon
  \Nat \to A$ दिया हो, तो $f(n) = g(n-1)$
  परिभाषित कीजिए।
\end{proof}

इससे हमें निम्नलिखित परिणाम मिलता है:

\begin{cor}\ollabel{cor:enum-nat}
समुच्चय $A$ !!{enumerable} है यदि और केवल यदि वह रिक्त है अथवा कोई
!!a{surjective} फलन $f\colon \Nat \to A$ है।
\end{cor}

हमने ऊपर चर्चा की कि समुच्चय~$A$ का हर !!{element}s लिखने वाली किसी सूची को
पुनरावृत्तियों-रहित सूची में बदला जा सकता है। गणन के लिए भी यह सत्य है, किंतु
इसे कठोरता से सूत्रबद्ध और सिद्ध करना कुछ कठिन है। प्रत्येक फलन
$f\colon \PosInt \to A$ को सभी $n \in
\PosInt$ के लिए परिभाषित होना चाहिए।
यदि~$A$ में केवल परिमित संख्या में !!{element}s हैं, तो स्पष्टतः हमारे पास
~$\PosInt$ पर परिभाषित ऐसा फलन नहीं हो सकता---उसके !!{element}s अनंततः अनेक
हैं---जो~$A$ का प्रत्येक !!{element}s मान के रूप में ले, पर कभी एक ही मान दो बार न ले।
उस स्थिति में, अर्थात जब पुनरावृत्तियों-रहित सूची परिमित हो, हमें~$f$ के लिए
भिन्न प्रांत चुनना होगा, जिसमें केवल परिमित संख्या में !!{element}s हों।
पुनरावृत्ति न होने का अर्थ है कि $f$ !!{injective} होना चाहिए। चूँकि वह
!!{surjective} भी है, हम किसी परिमित समुच्चय $\{1, \dots, n\}$ या
$\PosInt$ और~$A$ के बीच !!a{bijection} खोज रहे हैं।

\begin{prop}\ollabel{prop:enum-bij}
यदि $f\colon \PosInt \to A$ !!{surjective} है (अर्थात~$A$ का गणन है), तो
कोई !!a{bijection} $g\colon Z \to A$ है, जहाँ $Z$ या तो~$\PosInt$ है या
$\{1, \dots, n\}$ है, किसी~$n \in \PosInt$ के लिए।
\end{prop}

\begin{proof}
  हम फलन $g$ को पुनरावर्ती ढंग से परिभाषित करते हैं: $g(1) = f(1)$ मानें।
  यदि $g(i)$ पहले ही परिभाषित हो चुका है, तो $g(i+1)$ को $f(1)$, $f(2)$,
  \dots{} का ऐसा प्रथम मान लें जो पहले से $g(1)$, \dots, $g(i)$ में न हो,
  बशर्ते ऐसा कोई मान हो। यदि $A$ में ठीक $n$ !!{element}s हैं, तो $g(1)$,
  \dots, $g(n)$ सभी परिभाषित हैं; अतः हमने फलन $g\colon \{1, \dots, n\}
  \to A$ परिभाषित कर लिया है। यदि $A$ में अनंततः अनेक !!{element}s हैं, तो
  किसी भी $i$ के लिए~$A$ का कोई !!a{element} गणन $f(1)$, $f(2)$, \dots,
  में अवश्य होगा जो पहले से $g(1)$, \dots, $g(i)$ में नहीं है। इस
  स्थिति में हमने फलन $g\colon \PosInt \to A$ परिभाषित कर लिया है।
  
  फलन $g$ !!{surjective} है, क्योंकि~$A$ का प्रत्येक अवयव $f(1)$, $f(2)$,
  \dots{} में है ($f$ के !!{surjective} होने के कारण), और इसलिए अंततः~$g(i)$
  का मान होगा, किसी उपयुक्त~$i$ के लिए। वह !!{injective} भी है, क्योंकि यदि कोई
  $j < i$ ऐसा होता कि $g(j) = g(i)$, तो $g(i)$ पहले ही $g(1)$, \dots,
  $g(i-1)$ में होता, जो~$g$ की हमारी परिभाषा के विपरीत है।
\end{proof}

\begin{cor}\ollabel{cor:enum-nat-bij}
समुच्चय $A$ !!{enumerable} है यदि और केवल यदि वह रिक्त है अथवा कोई
!!a{bijection} $f\colon N \to A$ है, जहाँ या तो $N = \Nat$ है या
$N = \{0, \dots, n\}$ है, किसी $n \in \Nat$ के लिए।
\end{cor}

\begin{proof}
$A$ !!{enumerable} है यदि और केवल यदि $A$ रिक्त है अथवा कोई
!!a{surjective} फलन $f\colon \PosInt \to A$ है। \olref{prop:enum-bij} के
अनुसार, दूसरी बात तभी सत्य है जब कोई !!a{bijective} फलन~$f\colon Z \to A$
हो, जहाँ $Z =
\PosInt$ या $Z = \{1, \dots, n\}$ हो, किसी $n \in \PosInt$ के लिए।
\olref{prop:enum-shift} के प्रमाण वाली युक्ति
से, वह बात भी तभी सत्य है जब कोई !!a{bijection} $g\colon N \to A$ हो,
जहाँ या तो $N = \Nat$ या $N = \{0, \dots, n-1\}$ हो।
\end{proof}

\begin{prob}
  \olref[sfr][siz][enm]{defn:enumerable} के अनुसार, समुच्चय $A$ गणनीय है
  यदि और केवल यदि $A = \emptyset$ है अथवा कोई !!a{surjective} फलन $f\colon
  \PosInt \to A$ है। ``!!{enumerable} समुच्चय'' को इस प्रकार भी सटीकता से
  परिभाषित किया जा सकता है: कोई समुच्चय गणनीय है यदि और केवल यदि कोई
  !!a{injective} फलन $g\colon A \to \PosInt$ है। दिखाइए कि दोनों परिभाषाएँ
  समतुल्य हैं; अर्थात दिखाइए कि कोई !!a{injective} फलन
  $g\colon A \to \PosInt$ है यदि और केवल यदि या तो $A = \emptyset$ है या
  कोई !!a{surjective} फलन $f\colon \PosInt \to A$ है।
\end{prob}

\end{document}
